\scanpage{57}
% archive foliation: 56
% Sheet read after rotation through 180 degrees.

10.

5.2. \uline{Evaluation de la politique de recherche des établissements}

La D.E.S.U.R. ne sera crédible auprès de la D.G.R.S.T. et par conséquent n'obtiendra des crédits suffisants que si elle est capable de porter un jugement sur la valeur des résultats de recherche de chaque établissement et sur la façon dont les établissements conduisent leur politique de recherche. Actuellement les services de la D.E.S.U.R. disposent de nombreuses données sur les établissements :

\begin{itemize}
\item[-] le nombre des thèses soutenues
\item[-] le nombre et la qualité des publications
\item[-] le nombre d'associations avec les autres établissements de recherche
\item[-] le nombre des contrats
\item[-] évolution de la carrière du personnel
\end{itemize}

Le danger est grand de voir s'instituer un système purement numérique, alimenté par ces critères, et fournissant, par un ordinateur, \unsure{le:montant|le montant} des crédits de l'université. Juger la recherche conduite par un établissement est une opération délicate ; il est donc envisagé de placer, auprès du directeur des enseignements supérieurs et de la recherche un groupe de personnalités indépendantes capables de faire la synthèse des renseignements recueillis sur un établissement pour porter un jugement global qui restera confidentiel.

5.3. \uline{La responsabilité de la D.E.S.U.R.}

Les deux branches du Secrétariat d'Etat qui supportent la recherche dans les établissements d'enseignement supérieur, la D.E.S.U.R. et le C.N.R.S. doivent pratiquer une politique différente , quoique étroitement complémentaire. Le C.N.R.S. est une entreprise de recherche centralisée, qui agit par discipline ; il aide, dans un tel établissement tel laboratoire parce que cela est utile au programme de recherche. La D.E.S.U.R. agit par établissement ; elle aide chaque établissement à développer le programme qu'il propose, sous réserve qu'il contribue à l'une des missions spécifiques de la recherche universitaire.

Cela suppose la mise au point d'une politique de répartition des moyens complémentaires (\S\ 4.2.) qui en s'ajoutant aux moyens de base (\S4.1.) conditionnent la politique de recherche de l'établissement.

.../...

